% -------------------------------
% Paquetes y comandos añadidos para esta memoria.
% Aquí se añade lo que haga falta (TikZ, pgfplots, tablas...) en vez de tocar configuracion.tex.
% -------------------------------

% Utilidades para tablas y texto
\usepackage{booktabs}
\usepackage{etoolbox}
\usepackage{longtable}
\usepackage{csquotes}
\usepackage{icomma}
\usepackage{placeins}
\usepackage{xurl}
\usepackage{breakcites}  % permite partir citas largas entre líneas (evita overfull hbox con [Universidad de Castilla-La Mancha, 2026])

% Marcadores de claves usados en el diagrama entidad-relación (fig:modelo-datos).
% Se definen con \providecommand para no pisar nada si ya existieran.
\providecommand{\pk}{\textsc{PK}}
\providecommand{\fk}{\textsc{FK}}
\providecommand{\uk}{\textsc{UK}}
% \attr lista los atributos de una entidad en el diagrama ER.
% Acepta argumento opcional: \attr{nombre} o \attr nombre (el texto queda suelto).
\providecommand{\attr}[1][]{\texttt{#1}}

% Altura de cabecera recomendada por fancyhdr
\setlength{\headheight}{15pt}

% Permitir cortes en URLs en cualquier letra o dígito para evitar desbordes
\expandafter\def\expandafter\UrlBreaks\expandafter{\UrlBreaks%
  \do\a\do\b\do\c\do\d\do\e\do\f\do\g\do\h\do\i\do\j%
  \do\k\do\l\do\m\do\n\do\o\do\p\do\q\do\r\do\s\do\t%
  \do\u\do\v\do\w\do\x\do\y\do\z%
  \do\A\do\B\do\C\do\D\do\E\do\F\do\G\do\H\do\I\do\J%
  \do\K\do\L\do\M\do\N\do\O\do\P\do\Q\do\R\do\S\do\T%
  \do\U\do\V\do\W\do\X\do\Y\do\Z%
  \do\0\do\1\do\2\do\3\do\4\do\5\do\6\do\7\do\8\do\9\do\-}%

% Margen de tolerancia para justificación evitando overfull hbox
\setlength{\emergencystretch}{3em}

% Diagramas y gráficos
\usepackage{tikz}
\usetikzlibrary{babel,arrows.meta,positioning,shapes.geometric,fit,calc}

\usepackage{pgfplots}

% Diagramas de secuencia (sequencediagram, \newthread, \newinst, call, \mess)
\usepackage{pgf-umlsd}
\usepackage{adjustbox}  % max width para los dibujos TikZ
\pgfplotsset{compat=1.18}

% Enlaces sin recuadros de color
\hypersetup{hidelinks}

% Páginas en blanco (al empezar capítulo en página impar) sin cabecera ni número
\makeatletter
\renewcommand{\cleardoublepage}{\clearpage\if@twoside\ifodd\c@page\else
  \hbox{}\thispagestyle{empty}\newpage\if@twocolumn\hbox{}\newpage\fi\fi\fi}
\makeatother

% -------------------------------
% Maquetación de tablas y figuras (entrega Sprint 1)
% -------------------------------
% Celdas alineadas arriba (p) en vez de centradas en vertical (m): las filas ya no se pisan
\newcolumntype{L}[1]{>{\raggedright\let\newline\\\arraybackslash\hspace{0pt}}p{#1}}
\newcolumntype{C}[1]{>{\centering\let\newline\\\arraybackslash\hspace{0pt}}p{#1}}
\newcolumntype{R}[1]{>{\raggedleft\let\newline\\\arraybackslash\hspace{0pt}}p{#1}}
\renewcommand{\arraystretch}{1.15}
\AtBeginEnvironment{longtable}{\small}
\AtBeginEnvironment{table}{\small}
% Figuras y tablas: que compartan página con el texto en vez de quedarse solas
\renewcommand{\topfraction}{0.85}
\renewcommand{\bottomfraction}{0.6}
\renewcommand{\textfraction}{0.12}
\renewcommand{\floatpagefraction}{0.8}
\setcounter{topnumber}{3}
\setcounter{totalnumber}{4}
% Documento para leer en PDF: sin páginas en blanco entre capítulos
\renewcommand{\cleardoublepage}{\clearpage}
